\section{Variational Quantum Algorithms (VQAs)}\label{sec:variational-quantum-algorithms-vqas}

\subsection{Why Variational Algorithms?}\label{sec:why-variational-algorithms}

Most famous quantum algorithms assume quantum computers were perfect. For example: qubits remain coherent; gates are applied correctly; we can execute long circuits; we can use many qubits. However, on \autopageref{sec:why-quantum-computers-are-noisy} we discovered something unpleasant: \textbf{real quantum computers are noisy}. Qubits lose information because of decoherence, and gates are not perfect.

\highspace
\textcolor{Red2}{\faIcon{times-circle} \textbf{Can't we use error correction?}} In principle, we could use quantum error correction (\autopageref{sec:quantum-error-correction}) to mitigate noise. However there is a \textbf{problem}. A single logical qubit may require tens or hundreds of physical qubits to be protected against errors, and this could raise the \textbf{blow-up problem} (i.e., the number of physical qubits required to implement a quantum algorithm becomes unmanageable). So \hl{today's quantum computers do not yet have enough qubits to run large fault-tolerant computations.}

\highspace
\begin{flushleft}
    \textcolor{Green3}{\faIcon{book} \textbf{NISQ devices}}
\end{flushleft}
The term \definition{NISQ (Noisy Intermediate-Scale Quantum)} was introduced by John Preskill to \textbf{describe the current generation of quantum computers}. These machines are:
\begin{itemize}
    \item \important{Noisy}. Qubits are extremely fragile. They suffer from decoherence, gate errors, measurement errors, and environmental disturbances. The longer the circuit, the larger the error becomes.

    
    \item \important{Intermediate Scale}. Today's devices contain tens of qubits, hundreds of qubits are expected in the near future, and at most a few thousand physical qubits are expected in the next decade. This sounds impressive, but it is far from what is needed for large-scale fault-tolerant quantum computing.
\end{itemize}

\begin{definitionbox}[: NISQ Device]
    A \definition{NISQ (Noisy Intermediate-Scale Quantum) device} is a \textbf{current generation quantum computer} characterized by a limited number of qubits and significant noise sources such as decoherence and gate errors. Because these devices cannot reliably execute long fault-tolerant quantum algorithms, hybrid approaches such as VQE and QAOA were developed specifically for them.
\end{definitionbox}

\highspace
\begin{examplebox}[: An Analogy for the NISQ Problem]
    Imagine owning a Formula 1 engine. The engine can theoretically reach 350 km/h. Amazing. But suppose the fuel tank leaks after 10 seconds and we cannot finish the race. The engine is powerful, but the hardware is not mature enough. Current quantum computers are exactly like that:
    \begin{itemize}
        \item Quantum mechanics allows us to design powerful algorithms.
        \item Current hardware is not mature enough to run these algorithms without errors.
    \end{itemize}
\end{examplebox}

\highspace
\begin{flushleft}
    \textcolor{Green3}{\faIcon[regular]{lightbulb} \textbf{The Key Idea: Hybrid Quantum-Classical Approaches}}
\end{flushleft}
Researchers asked: ``\emph{instead of waiting for perfect quantum computers, can we already do something useful with today's noisy machines?}'' This is the birth of \textbf{variational quantum algorithms (VQAs)}, which are designed to be robust against noise and to work within the limitations of NISQ devices.

\highspace
The core idea behind variational quantum algorithms is to \textbf{combine quantum and classical computing resources} in a hybrid approach. The quantum computer is used for a specific part of the computation, while classical methods wrap this quantum subroutine and optimize it.

\highspace
\textcolor{Green3}{\faIcon{question-circle} \textbf{Why split the computation?}} Because the \textbf{two machines have different strengths}. Classical computers are very good at optimization, arithmetic, memory management, parameter updates, and decision making. Instead, quantum computers are potentially very good at representing large quantum states, exploiting superposition, exploiting entanglement, and estimating quantum properties.

\highspace
\textcolor{Green3}{\faIcon{question-circle} \textbf{What does the Classical Computer do?}} The \textbf{classical computer acts like a coach}. Suppose the quantum circuit has a parameter $\theta$. The quantum computer runs and returns a cost function $C(\theta)$. The classical computer takes this cost, updates the parameter $\theta$ (e.g., using gradient descent), and sends the new parameter back to the quantum computer. This process is repeated until the optimizer (i.e., the classical computer) finds the optimal parameter $\theta^*$.

\highspace
In this way, the quantum computer is not executing a long, complex algorithm on its own. Instead, it is performing many short executions of a parameterized quantum circuit, which is more feasible on NISQ devices. The classical computer handles the optimization and decision-making, while the quantum computer provides the quantum advantage where it can.

\highspace
\begin{examplebox}[: Analogy]
    Think about a Formula 1 race. \emph{Who wins the race?} Not only the driver, but also the engineers, the strategists, and the pit crew (i.e., the whole team). The driver does what only the driver can do, while the engineers do everything else.
\end{examplebox}

\highspace
\begin{examplebox}[: Optimizing the Quantum Circuit, Not the Solution]
    Imagine we want to find the minimum of a function $f(x)$. Classically we would directly try different values of $x$. However, in a Variational Algorithm (VQA), we do something slightly different. The parameter $\theta$ controls a quantum circuit and the optimizer changes $\theta$ until the cost function becomes minimal. So we are not directly optimizing the solution. We are optimizin the \textbf{quantum circuit that generates the solution}.
\end{examplebox}